\section{From a volatility smile to a probability law}
\label{sec:intro}

A listed option board contains finitely many strikes.  A pricing system needs
a curve at every strike, including strikes that are not quoted.  Traders view
that curve as a volatility smile.  Mathematically it is the option-price image
of a probability distribution.  The two views are equivalent, but the second
one makes the restrictions clear: option prices must decrease with strike,
butterflies must have nonnegative value, the terminal asset price must remain
positive, and its mean must equal the forward.

The LQD model starts from this distribution rather than from volatility.  Its
central device is a monotone map from percentile rank to log-return.  A
positive speed makes the map increasing and therefore makes the density
positive.  The smile is calculated only after the law has been built.

The conceptual core is contained in
\cref{sec:intro,sec:model,sec:validity,sec:ledger}: law, transport, tails, and
pricing.  \Cref{sec:calibration,sec:numerics,sec:calendar} turn those identities
into a fitted surface.  The
proofs and implementation details are collected in the chapter appendices so
that the main argument can be read without interrupting its progression.

\subsection{Normalize first}

Fix an expiry $T$.  Let $F$ be its forward and $D$ its discount factor.  Work
under the $T$-forward measure and set
\begin{equation}
X=\log\frac{S_T}{F},\qquad Y=e^X=\frac{S_T}{F},\qquad \E[Y]=1.
\label{eq:Xdef}
\end{equation}
The expectation may simply be read as the normalized pricing expectation; no
dynamic model is needed in this chapter.  The last equality is the forward,
or martingale, condition.  For a strike $K$ write
\[
k=\log(K/F),\qquad y=e^k=K/F.
\]
Writing $a^+=\max(a,0)$, the normalized call and put prices are
\begin{equation}
c(k)=\E[(Y-e^k)^+],\qquad
P(k)=\E[(e^k-Y)^+],\qquad
c(k)-P(k)=1-e^k.
\label{eq:calldef}
\end{equation}
Multiplying by $DF$ returns the cash price.  This normalization removes rates,
dividends, and currency units from the one-expiry problem.

Markets usually quote implied volatility rather than normalized price.  If
$\tau$ is the chosen year fraction and $w(k)=\sigma(k)^2\tau$ is total implied
variance, the normalized Black call is
\begin{equation}
\Black(k,w)=\PhiN(d_+)-e^k\PhiN(d_-),
\qquad d_\pm=-\frac{k}{\sqrt w}\pm\frac{\sqrt w}{2}.
\label{eq:black}
\end{equation}
For every admissible price strictly between intrinsic value and the forward,
the implied variance $w(k)$ is the unique positive solution of
$\Black(k,w(k))=c(k)$ \citep{Black1976}.  Price and implied volatility are
therefore two coordinates for the same object.

\subsection{The call curve and the law}

Use normalized strike $y$, not log-strike $k$, for the next calculation.  If
$Y$ has density $f_Y$, differentiating the call payoff gives
\begin{equation}
\frac{\partial c}{\partial y}=-\Prob(Y>y),\qquad
\frac{\partial^2c}{\partial y^2}=f_Y(y)\ge0.
\label{eq:BL}
\end{equation}
This is the Breeden--Litzenberger identity
\citep{BreedenLitzenberger1978}.  The slope of the call curve is minus the
probability of finishing above strike; its curvature is the terminal density.

The same identity has a traded reading.  For a strike spacing $\delta$, the
symmetric butterfly
\[
c(y-\delta)-2c(y)+c(y+\delta)
\]
has a nonnegative payoff in every terminal state.  Dividing its price by
$\delta^2$ and shrinking $\delta$ recovers $f_Y(y)$.  Negative curvature
would give a negative price to a nonnegative payoff.

\begin{definition}[Arbitrage-free slice]
\label{def:slice}
An arbitrage-free slice is a call curve generated by a positive random
variable $Y$ with $\E[Y]=1$.  If $Y$ has no atom at zero and has a density,
the curve is equivalently a decreasing convex function of $y>0$ satisfying
\[
(1-y)^+\le c(y)\le1,\qquad -1\le c'(y)\le0,
\qquad c(y)\longrightarrow0\quad(y\to\infty),
\]
with $c'(y)\to-1$ as $y\downarrow0$.
\end{definition}

Every condition has a direct source.  The upper bound follows from
$(Y-y)^+\le Y$; the lower bound is Jensen's inequality; the slope is a
survival probability; and the right boundary follows from integrability of
$Y$.  Conversely, the curvature of such a curve defines a probability
measure with mean one.  Thus a valid call curve and a positive mean-one law
contain exactly the same information.

One distinction matters in practice: convexity belongs to strike $y$, not to
log-strike $k$.  The chain rule gives
\[
\frac{\partial^2c}{\partial k^2}
=e^k\{e^kf_Y(e^k)-\Prob(Y>e^k)\},
\]
which is negative far in the left wing for an ordinary density.  Testing
convexity in log-strike rejects correct curves.

In volatility coordinates the same requirement becomes less transparent.
Twice differentiating $c(k)=\Black(k,w(k))$ yields
\begin{equation}
f_X(k)=g_{\rm D}(k)\frac{\phin(d_-(k,w(k)))}{\sqrt{w(k)}},
\quad
g_{\rm D}(k)=
\left(1-\frac{k w'}{2w}\right)^2
-\frac{(w')^2}{4}\left(\frac1w+\frac14\right)+\frac{w''}{2}.
\label{eq:durrleman}
\end{equation}
A twice differentiable smile is butterfly-free exactly when
$g_{\rm D}(k)\ge0$ for every $k$ \citep{GatheralJacquier2014}; the calculation
is given in \cref{app:durrlemanproof}.  Direct smile models must control this
nonlinear expression.  LQD instead constructs $f_X>0$, so the condition is a
consequence rather than a calibration constraint.

\subsection{Four coordinates, used for four different jobs}

The chapter changes horizontal coordinate several times.  Each change has a
specific purpose:
\[
\begin{array}{c|c|c}
\text{coordinate}&\text{meaning}&\text{main use}\\ \hline
y=K/F&\text{normalized strike}&\text{convexity and static arbitrage}\\
k=\log y&\text{log-moneyness}&\text{market smiles}\\
u\in(0,1)&\text{percentile rank}&\text{probability and quantiles}\\
z=\log\{u/(1-u)\}&\text{log-odds}&\text{integration and numerics}
\end{array}
\]
The logit map $u\mapsto z$ is a bijection from $(0,1)$ to $\R$.  Its inverse
and the logistic density are
\begin{equation}
u=\Lambda(z)=\frac1{1+e^{-z}},\qquad
\rho(z)=\Lambda(z)\Lambda(-z).
\label{eq:logodds}
\end{equation}
The median is $z=0$; the tenth and ninetieth percentiles are
$z=\mp\log9$.  Moving one unit in $z$ multiplies the odds by $e$.  This turns
the open rank interval into an ordinary real line and makes both tails
numerically accessible.

The complete construction follows one direction:
\[
\text{rank}\ u
\longrightarrow \text{log-odds}\ z
\longrightarrow \text{return}\ x(z)
\longrightarrow \text{law}\ f_X
\longrightarrow \text{price}\ c
\longrightarrow \text{smile}\ w.
\]
Rank supplies exactly one unit of probability.  Positive transport speed
makes the return map increasing.  The Black formula is used only at the final
change of coordinates.  This order is the organizing principle of the
chapter.
